\section{System Design}
\label{sec:overview}

This section describes who takes part, what each party is trusted with, and what happens to one request. Three choices
shape the design. Everything that touches tenant keys, or makes scheduling decisions, runs inside the attested image, so
the rest of the system only has to cope with machines that stop or answer late (Section~\ref{sec:trust}). Trust, keys and
warm sandboxes are prepared ahead of demand, because they take seconds to minutes to establish
(Section~\ref{sec:sched}). And a machine is judged by whether the work it accepted comes back, as seen by the router or the
scheduler, not by what the machine says about itself (Section~\ref{sec:ft}).

\subsection{What the provider stores, and why a compromised provider learns little}
\label{sec:atrest}
A tenant encrypts its code bundle with its own key before uploading it. The provider's registry, the disks of workers that
cache the bundle, and any snapshot of a running sandbox hold only this ciphertext; snapshots are encrypted with the same
key when they are written. The key itself never passes through the provider: it travels from the tenant's key service
directly to an attested worker, over a channel that ends inside the attested image, and stays in that worker's memory until
the session ends. As a result, an attacker who takes over the provider's registry, database, router or broker obtains
encrypted bundles and request metadata, but no tenant code and no tenant data.

\subsection{The parties}
Figure~\ref{fig:arch} shows the components for the clustered architecture of Section~\ref{sec:arch2}.

\begin{figure}[t]
\centering
% Architecture figure (Fig. 1). Native width 139 mm (acmart single column). y runs downwards.
% Rows: clients and legend; provider control plane; volunteer cluster; tenant domain.
\begin{tikzpicture}[x=1mm,y=1mm]
  \tikzset{box/.append style={font=\footnotesize}}
  % ---------- clients and legend ----------
  \node[box, minimum width=20mm, minimum height=7mm] (client) at (52,-3.5) {Clients};
  \begin{scope}[font=\tinyx, every node/.style={inner sep=1pt}]
    \node[prov, minimum size=3mm] (l1) at (75,-3.5)  {};  \node[anchor=west] at (l1.east) {provider};
    \node[att,  minimum size=3mm] (l2) at (92,-3.5)  {};  \node[anchor=west] at (l2.east) {attested code};
    \node[ten,  minimum size=3mm] (l3) at (114,-3.5) {};  \node[anchor=west] at (l3.east) {tenant};
    \node[untr, minimum size=3mm] (l4) at (128,-3.5) {};  \node[anchor=west] at (l4.east) {untrusted};
  \end{scope}
  % ---------- provider control plane ----------
  \node[zone, fill=blue!4, minimum width=139mm, minimum height=29mm, anchor=north west] (prov) at (0,-10) {};
  \node[font=\scriptsize\bfseries, anchor=north west] at (2,-10.8) {Provider control plane};
  \node[font=\tinyx, anchor=north east] at (137,-11.2) {availability and accounting, no tenant keys};
  \node[prov, minimum width=31mm, minimum height=15mm, font=\scriptsize] (ver)  at (17,-29)  {\textbf{Verifier, registry}\\[-1pt]\tinyx session certificates\\[-2pt]\tinyx encrypted bundles};
  \node[prov, minimum width=31mm, minimum height=15mm, font=\scriptsize] (rtr)  at (52,-29)  {\textbf{Routers}\\[-1pt]\tinyx pending copies\\[-2pt]\tinyx progress watchdogs\\[-2pt]\tinyx receipts};
  \node[prov, minimum width=31mm, minimum height=15mm, font=\scriptsize] (brk)  at (87,-29)  {\textbf{Broker}\\[-1pt]\tinyx one topic set\\[-2pt]\tinyx per cluster};
  \node[prov, minimum width=31mm, minimum height=15mm, font=\scriptsize] (orch) at (122,-29) {\textbf{Orchestrator}\\[-1pt]\tinyx clusters, epochs\\[-2pt]\tinyx database};
  % ---------- volunteer cluster ----------
  \node[zone, fill=gray!5, minimum width=139mm, minimum height=44mm, anchor=north west] (vol) at (0,-48) {};
  \node[font=\scriptsize\bfseries, anchor=north west] at (2,-48.8) {Volunteer cluster: tens of machines};
  \node[att, minimum width=84mm, minimum height=21mm, font=\scriptsize] (w) at (45,-69)
       {\textbf{Worker (attested appliance)}\\[2pt]{\tinyx\begin{tabular}{@{}l@{\hspace{4mm}}l@{}}agent and TPM quote & keys in memory only\\ Firecracker microVMs & sealed, encrypted cache\end{tabular}}};
  \node[att, minimum width=38mm, minimum height=21mm, font=\scriptsize] (sch) at (115,-69)
       {\textbf{Worker with}\\[-1pt]\textbf{scheduler role}\\[-1pt]\tinyx provisioning and request placement\\[-2pt]\tinyx rebuildable state, epoch $e$};
  \node[untr, minimum width=73mm, minimum height=6mm, font=\tinyx] (own) at (99,-87) {Owner controls hardware, power, network: untrusted};
  % ---------- tenant domain ----------
  \node[zone, fill=orange!6, minimum width=66mm, minimum height=21mm, anchor=north west] (tenant) at (3,-96) {};
  \node[font=\scriptsize\bfseries, anchor=north west] at (5,-96.8) {Tenant domain};
  \node[ten, minimum width=58mm, minimum height=8mm, font=\scriptsize] (kbs) at (36,-110) {\textbf{Key service}: secrets to attested hosts only};
  % ---------- numbered flows (one invocation); dashed = control traffic off the hot path ----------
  \draw[msg] (client.south) -- node[lbl, left=0.5mm] {\circled{1}} (rtr.north);
  \draw[msg] (rtr.east) -- node[lbl, above=2.5mm] {\circled{2}} (brk.west);
  \draw[msg, <->] (99,-36.5) -- node[lbl, right=0.5mm, pos=0.3] {\circled{2}\circled{3}} (99,-58.5);
  \draw[msg] (82,-36.5) -- node[lbl, right=0.5mm, pos=0.3] {\circled{3}} (82,-58.5);
  \draw[<->, thick] (52,-36.5) -- node[lbl, right=0.5mm, align=left, pos=0.3] {\circled{4} input, result, receipt} (52,-58.5);
  \draw[msg, <->] (11,-36.5) -- node[lbl, left=0.3mm, pos=0.3] {\circled{5}} (11,-58.5);
  \draw[msg, <-] (23,-58.5) -- node[lbl, right=0.5mm, pos=0.3] {\circled{7}} (23,-36.5);
  \draw[msg, dashed, <->] (brk.east) -- node[lbl, above=2.5mm] {\circled{8}} (orch.west);
  \draw[msg, <->] (36,-79.5) -- node[lbl, right=0.5mm] {\circled{6} key request} (36,-96);
  \node[zone, dashed, fill=gray!4, minimum width=58mm, minimum height=21mm, anchor=north west, font=\tinyx, align=center] (other) at (77,-96) {Other clusters have the same structure\\and their own topics on the same broker};
\end{tikzpicture}

\caption{The \voss{} architecture with geographic clusters. Numbers follow one invocation: (1) a client calls a router; (2) the router sends the request's metadata to the cluster's scheduler through the broker; (3) the scheduler sends the assignment to a worker; (4) the router sends the input and the worker returns the result and a signed receipt; (5) the worker proves itself to the verifier; (6) it asks the tenant's key service for the key; (7) it fetches the encrypted code from the registry. (8) The orchestrator appoints schedulers, off the request path.}
\label{fig:arch}
\end{figure}

\textbf{Provider services} are assumed reliable and are trusted to keep running and to keep correct records, but never with
tenant keys.
The \emph{router} is the public entry point. A client sends its request to the nearest router. The router chooses a
cluster, sends the request's metadata (not its input) to that cluster's scheduler, and keeps its own copy of the request
until a worker confirms it has started. If the scheduler does not answer within 1\,s, the router resends the request or
sends it to another cluster. When the worker is chosen, the router streams the input to it, receives the result and returns
it to the client, and checks and stores the receipt the worker signs for that work.
The \emph{orchestrator} decides which workers form each cluster and appoints each cluster's scheduler.
The \emph{verifier} checks attestation evidence and acts as the certificate authority: it issues each attested worker a
session certificate and revokes it when the session ends (Section~\ref{sec:trust}).
The \emph{registry} stores encrypted code bundles and file systems, and the \emph{broker} carries messages between a
cluster's machines (Section~\ref{sec:network}).

\textbf{Tenants} run, or entrust to a service they choose, a \emph{key service} that holds their keys and releases one only
to a worker the verifier has accepted. The provider does not run it.

\textbf{Volunteer machines} reboot into the \emph{attested image}: a locked, measured operating system containing the worker
software, the Firecracker microVM monitor~\cite{agache2020firecracker} and the scheduler code. Each worker runs every
invocation in its own microVM and keeps an encrypted cache of code and snapshots. While the image runs the owner cannot use
the machine, so donation means a scheduled or dedicated session rather than background use of a desktop. The owner still
controls power, network and firmware settings, which is why a worker can stop or answer late even though its code cannot
misbehave.

\subsection{Two architectures}
\label{sec:arch2}
We designed and compared two ways to organise scheduling.

\emph{One global scheduler.} A single scheduler, run on provider hardware, plans and places for every worker. Workers open
connections to it at a stable address, and its failures are handled by ordinary replication inside the provider. This is
simple, and it lets the planner pool all machines when deciding where keys and warm sandboxes go.

\emph{Geographic clusters.} The orchestrator groups workers into clusters of tens of machines by network distance, and each
cluster has its own scheduler. Each planning problem stays small and the scheduler is close to its workers. The scheduler
runs on one of the cluster's own volunteer machines: the scheduler code is part of the image, so any worker can take the
role without booting anything, and the provider runs no per-cluster server. The price is that a scheduler can now fail like
any volunteer machine, which Section~\ref{sec:ft} handles, and that it is hard to reach, which Section~\ref{sec:network}
handles with a publish/subscribe broker.

Both use the same planner, per-request step and workers. Section~\ref{sec:res-nopay} compares them: at the scale we can
simulate, the global scheduler was faster.

\subsection{Life of an invocation}
The router sends the request's metadata to the chosen cluster's scheduler, which picks a worker and sometimes a second one for a duplicate copy (a \emph{hedge}). The router then sends the input to the worker over a connection that ends inside the attested image, so the scheduler never sees inputs, results or keys. A worker without the tenant's key first requests it from the key service, then starts the function at whatever level it has reached and runs it. The result returns to the router with a signed receipt (Section~\ref{sec:receipts}). Each part of the latency is logged separately, which lets Section~\ref{sec:results} attribute delays to their causes.

\subsection{Sessions}
\label{sec:lifecycle}
A worker becomes usable only after booting and passing attestation, which takes \CalAttested\,s from power-on on our laptop. When the
session ends, with or without notice, every key and warm sandbox on the worker is gone. Cached bundles and snapshots stay on
disk, still encrypted, and are usable again after the next attestation and key release. So within a session each function
on a worker is at one of the levels L0 to L3 of Section~\ref{sec:background}, and separately the worker either holds or
does not hold the tenant's key. The simulator reads these two facts, and their costs, from a data file for each trust
technology (Section~\ref{sec:substrates}).

\subsection{Rules for messages and state}
Machines interact only by typed messages, and every connection is mutual TLS: both ends present certificates, and a worker
or scheduler presents the session certificate the verifier issued after attestation. The provider's database is private
to the provider services and is never on the request path, because a database round trip would add to every request's
latency. State lives where it is created: queues, caches and keys on workers, unassigned requests at routers, estimates in
the scheduler. After a scheduler is replaced, its successor rebuilds its view by asking the workers
(Section~\ref{sec:succession}).
